\begin{abstracteng}

Emission spectra consist of light emitted by electrons transitioning between the energy levels of atoms or ions. The unique energy level structure of each atom or ion forms a distinct pattern, enabling precise analysis of material composition using the \textit{Laser-Induced Breakdown Spectroscopy} (LIBS) method. However, the accuracy of this conventional method is often hindered by the complexity of multi-element emission spectra, matrix effects, and plasma fluctuations. This research develops a \textit{Deep Learning} (DL) based analysis approach that learns from simulated synthetic emission spectra. These synthetic spectra are simulated using the Saha-Boltzmann equations, where each emission line is broadened using a Gaussian line profile under the assumption of \textit{Local Thermal Equilibrium} (LTE), utilizing atomic transition data from the NIST \textit{Atomic Spectra Database} (ASD). The Informer model, a Transformer architecture optimized with \textit{ProbSparse Self-Attention}, \textit{Self-Attention Distilling}, and a \textit{Generative-style Decoder}, was chosen to handle the long sequences of emission spectra. The model is trained and evaluated using the F1-Score metric to achieve an optimal balance between precision and recall for element identification. This approach provides a strong proof of concept for reliable multi-element detection in the LIBS method and contributes to the advancement of modern spectroscopy.

\bigskip
\noindent
\textbf{Keywords:} LIBS, Informer, Synthetic Emission Spectra, Deep Learning, Predictive Analysis, ProbSparse Self-Attention.
\end{abstracteng}
